\section{The wall at rank two}\label{sec:wall}

The family's range is $m\in(1,2]$, the first gap of the rank sequence. This
section shows that its upper endpoint is not a convention but the zero of
$a=\frac{2-m}2$, and that three distinct things happen there. Two of them are in
this section; the third, which is what the endpoint costs, is
Section~\ref{sec:criterion}.

\subsection{The factorization, and where the two poles come from}\label{sec:poles}

\begin{proposition}[Factorization {\cite[Prop.~3.1]{Markov}}]\label{prop:factor}
$K_\omega=R_\omega\Kt_\omega$ with
\begin{equation}\label{eq:factor}
 R_\omega(p)=\frac{(s-\omega)(s-\omega-1)}{(s+\omega)(s+\omega-1)}
 =\frac{(p+a)(p-b)}{(p+b)(p-a)}
 =1-\frac{4\omega b}{p+b}-\frac{4\omega a}{p-a} .
\end{equation}
Apart from its poles at the zeros of $\Lambda(p+b)$ and its zeros at the zeros of
$\Lambda(p+a)$, $\Kt_\omega$ has exactly two poles, at $p=b$ and $p=-a=c$, and
exactly two zeros, at $p=a$ and $p=-b$; $R_\omega$ cancels all four. Both factors
are unimodular on $\re p=0$.
\end{proposition}

The two cancelled poles are the poles of $\Lambda(p+a)$, at $p+a=1$ and $p+a=0$.
They have different origins, and the distinction is what makes the repair of
Section~\ref{sec:cmpast} the natural one.

\begin{table}[ht]\centering\small
\caption{The two poles of $\Kt_\omega$ in the right half-plane, by origin.
Writing $\Lambda(p+a)=\pi^{-(p+a)/2}\Gamma(\frac{p+a}2)\zeta(p+a)$, the pole at
$p=b$ comes from $\zeta$ and the pole at $p=c$ from $\Gamma$. The remaining poles
of $\Gamma(\frac{p+a}2)$, at $p+a=-2,-4,\dots$, are cancelled by the trivial
zeros of $\zeta(p+a)$, and $\zeta(0)=-\frac12\ne0$, so the pole at $p+a=0$ is
exactly the archimedean one.}\label{tab:poles}
\begin{tabular}{@{}lll@{}}
\toprule
pole & origin & in $\re p>0$ iff\\
\midrule
$p=b=\frac m2$ & $\zeta(p+a)$ at $p+a=1$: the \emph{comb} & always\\
$p=c=\frac{m-2}2$ & $\Gamma(\frac{p+a}2)$ at $p+a=0$: the \emph{archimedean factor} & $m>2$\\
\bottomrule
\end{tabular}
\end{table}

\begin{proposition}[The second factor is a Blaschke factor exactly past the
wall]\label{prop:twoblaschke}
For every $\omega>0$,
\begin{equation}\label{eq:twoblaschke}
 R_\omega(p)=B_b(p)\cdot\frac{p+a}{p-a},\qquad\text{and}\qquad
 \frac{p+a}{p-a}=B_c(p)\ \text{ with }c=\omega-\tfrac12 .
\end{equation}
Consequently
\begin{itemize}
\item for $\omega<\frac12$ one has $c<0$, and $R_\omega=B_b/B_a$ has a pole at
 $p=a>0$: it is unimodular on the axis but \emph{not} inner, and
 $\sup_{\re p>0}\lvert R_\omega\rvert=\infty$;
\item for $\omega\ge\frac12$ one has $c\ge0$, and $R_\omega=B_bB_c$ is a product
 of two Blaschke factors of the right half-plane, hence \emph{inner}:
 $\lvert R_\omega\rvert\le1$ on $\re p>0$.
\end{itemize}
At $\omega=\frac12$ exactly, $c=0$ and $B_c\equiv1$: the two regimes meet at the
single factor $B_b$.
\end{proposition}

\begin{proof}
Equation~\eqref{eq:factor} and $a-1=-b$, $b-1=-a$ give the first identity; and
$\frac{p+a}{p-a}=\frac{p-(-a)}{p+(-a)}=B_{-a}=B_c$. A Blaschke factor with
positive parameter is inner and a product of inner functions is inner; for
$\omega<\frac12$ the pole at $p=a>0$ lies in the open right half-plane.
\end{proof}

\subsection{Complete monotonicity survives the crossing}\label{sec:cmpast}

The parent manuscript proves that $\Kt_\omega$ is completely monotone on
$(b,\infty)$ \cite[Prop.~3.2]{Markov} and that
$\Kh_\omega=\frac{p+a}{p-a}\Kt_\omega$ is too \cite[Prop.~3.6]{Markov}, so that
$K_\omega=B_b\Kh_\omega$ exhibits the transfer as a positive measure minus twice
its exponential moving average. The first of these carries no restriction on
$\omega$; the second, as proved there, does. We record both facts and repair the
second.

\begin{proposition}[Complete monotonicity, uncapped]\label{prop:cmuncapped}
For every real $\omega>0$:
\begin{enumerate}
\item[(a)] $\Kt_\omega$ is completely monotone on $(b,\infty)$;
\item[(b)] $\Kh_\omega(p)=\frac{p+a}{p-a}\Kt_\omega(p)$ is completely monotone on
 $(b,\infty)$.
\end{enumerate}
For $\omega\le\frac12$ part (b) is \cite[Prop.~3.6]{Markov}. For
$\omega\ge\frac12$ its proof there fails and the following grouping replaces it:
with $z=\frac{p+a}2$ and $\Kh_\omega=\Kh^\Gamma_\omega K^\zeta_\omega$,
\begin{equation}\label{eq:regroup}
 \Kh^\Gamma_\omega(p)
 =\pi^\omega\,\frac{\Gamma(z+c)}{\Gamma(z+\omega)}\cdot
 \frac{\Gamma(z+1)}{\Gamma(z+c+1)} ,
\end{equation}
a product of two Gamma ratios whose parameter gaps are $\frac12$ and $c$.
\end{proposition}

\begin{proof}
(a) is \cite[Prop.~3.2]{Markov}; its proof uses only that $p>b$ implies
$p+a>1$, which holds for every $\omega>0$, and is therefore uncapped.

For (b) with $\omega\le\frac12$, the parent writes
$\frac{p+a}{p-a}=1+\frac{2a}{p-a}$, whose causal kernel is
$\delta_0+2a\,e^{at}\mathbf 1_{t>0}$; this is a positive measure if and only if
$a\ge0$, that is if and only if $\omega\le\frac12$. Past the wall the same
expression has $2a<0$ and the argument is unavailable.

For $\omega\ge\frac12$ regroup instead. From
$(p+a)\Gamma\!\big(\frac{p+a}2\big)=2\Gamma(z+1)$, from $\frac{p+b}2=z+\omega$,
and from $p-a=2(z+c)$,
\[
 \frac{p+a}{p-a}\cdot\pi^\omega\frac{\Gamma(\frac{p+a}2)}{\Gamma(\frac{p+b}2)}
 =\pi^\omega\,\frac{\Gamma(z+1)}{(z+c)\,\Gamma(z+\omega)}
 =\pi^\omega\,\frac{\Gamma(z+c)}{\Gamma(z+\omega)}\cdot
 \frac{\Gamma(z+1)}{\Gamma(z+c+1)},
\]
the last step by $\Gamma(z+c+1)=(z+c)\Gamma(z+c)$. The parameter gaps are
$\omega-c=\frac12>0$ and $(c+1)-1=c\ge0$, so each factor is completely monotone
in $z$ by \eqref{eq:betadensity} (the second is identically $1$ when $c=0$), and
$z=\frac{p+a}2$ is an increasing affine function of $p$ with $z>\frac12$ on
$p>b$, inside both domains. The comb $K^\zeta_\omega$ is completely monotone on
$(b,\infty)$ with the positive coefficients \eqref{eq:comb}. Products of
completely monotone functions are completely monotone.
\end{proof}

\begin{remark}[The two groupings are exchanged at the wall]\label{rem:exchange}
Proposition~\ref{prop:cmuncapped} proves the same statement twice, by arguments
that are available on complementary ranges and meet where both degenerate. At
$\omega=\frac12$ the parent's factor $\frac{p+a}{p-a}$ is identically $1$ and so
is the second Gamma ratio of \eqref{eq:regroup}. At $\omega=1$ the two Gamma
ratios telescope: $c=\frac12$, $\omega=1$, and
$\Kh^\Gamma_1(p)=\pi\big/(z+\frac12)=2\pi/(p+\frac12)$, whose kernel is
$2\pi e^{-\tau/2}$.
\end{remark}

\begin{remark}[The trade]\label{rem:trade}
Table~\ref{tab:trade} collects what changes at $m=2$ and what does not. The
factor $\frac{p+a}{p-a}$ does not \emph{lose} a positivity there: it
\emph{exchanges} complete monotonicity for innerness. Everything else in the
decomposition is positive on both sides. Whether that exchange is the same fact
as the failure of the lattice on the same interval is the question of
Section~\ref{sec:position}, and the answer is no.
\end{remark}

\begin{table}[ht]\centering\small
\caption{What the wall separates. Only the first row changes, and it changes by
exchange rather than by loss.}\label{tab:trade}
\begin{tabular}{@{}lll@{}}
\toprule
object & $m<2$ & $m>2$\\
\midrule
$\frac{p+a}{p-a}$ & completely monotone, \emph{not} inner & inner, \emph{not} completely monotone\\
$R_\omega=B_b\cdot\frac{p+a}{p-a}$ & not inner (pole at $p=a>0$) & inner\\
$\Kt_\omega$ & completely monotone & completely monotone\\
$\Kh_\omega$ & completely monotone (parent's grouping) & completely monotone (regrouped)\\
$k^\Gamma_\omega,\ \ct_n$ & positive & positive\\
\bottomrule
\end{tabular}
\end{table}
